% Regression: t2_mpo — renders the 4-valent MPO tensor — red virtual MPO wire a: j -> k.
% Target: RMP 2011.12127 fig3_mpo.pdf (the MPO tensor as a fat graph).
% Formula: the 4-valent MPO tensor — red virtual MPO wire  a: j -> k
% (horizontal, directed east) crossing the black physical wire
% alpha: n -> m (vertical, directed north); fat-graph faces A, B, C, D
% in the four quadrants, gray wedges at the wire endpoints j, k, m, n.
%
% The OBJECT is grid-expressible (one op-row MPO cell), so panel 1 states
% it in the grid language; the fat-graph DEPICTION is genuinely non-grid,
% so panel 2 attempts it in the free tier.
% Honest gaps (free panel): no free-tier annotation-label primitive, so
% the quadrant faces A, B, C, D cannot be labeled; \tnjoin labels sit
% fixed north of the chord midpoint, so `a` collides with the crossing
% and a mid-wire alpha on the vertical join is unplaceable; no wire-hue
% key (red vs black lost); no gray-wedge terminal skin.
\documentclass[varwidth=19cm, border=6pt]{standalone}
\usepackage{amsmath, amssymb}
\usepackage{tenkz}
\begin{document}

% panel 1 — the same tensor in the grid language: an MPO site with open
% virtual indices j (west), k (east) and physical legs m (up), n (down);
% bond dir=right orients the virtual level j -> k as in the paper.
\[
  \begin{tenkz}[physical=updown, west label=$j$, east label=$k$,
                bond dir=right]
    \tn[up=$m$, down=$n$]{}
  \end{tenkz}
\]

% panel 2 — the fat-graph depiction, free tier: the two directed wires
% crossing with no glyph at the crossing (the fat graph IS the tensor).
\[
  \begin{tenkzfree}
    \tnput[dot, label pos=west,  ports={east:virtual}]{j}{(0,0)}{j}
    \tnput[dot, label pos=east,  ports={west:virtual}]{k}{(26mm,0)}{k}
    \tnput[dot, label pos=south, ports={north:physical}]{n}{(13mm,-13mm)}{n}
    \tnput[dot, label pos=north, ports={south:physical}]{m}{(13mm,13mm)}{m}
    \tnjoin[dir, label=a]{j.east}{k.west}
    \tnjoin[dir]{n.north}{m.south}
  \end{tenkzfree}
\]

\end{document}
